\corrigesChapitre{12}{L’énergie : transferts et conversions}
\par\noindent\begin{minipage}{\linewidth}
\begin{corrige}{1}{Identifier une source}
\begin{questions}
\item La pile, via le circuit.
\item De l’énergie chimique.
\end{questions}
\end{corrige}
\end{minipage}
\par\noindent\begin{minipage}{\linewidth}
\begin{corrige}{2}{Nommer une forme}
\begin{questions}
\item L’énergie cinétique.
\item Une énergie potentielle liée à sa position dans le système Terre-objet.
\end{questions}
\end{corrige}
\end{minipage}
\par\noindent\begin{minipage}{\linewidth}
\begin{corrige}{3}{Une unité}
\begin{questions}
\item Le joule, J.
\item Non : le watt est une unité de puissance, étudiée au chapitre 17.
\end{questions}
\end{corrige}
\end{minipage}
\par\noindent\begin{minipage}{\linewidth}
\begin{corrige}{4}{Une sortie moins visible}
\begin{questions}
\item Un transfert par rayonnement et des transferts thermiques.
\item Non : l’énergie est transférée à l’environnement.
\end{questions}
\end{corrige}
\end{minipage}
\par\noindent\begin{minipage}{\linewidth}
\begin{corrige}{5}{Renouvelable ou fossile}
\begin{questions}
\item Le vent relève d’une ressource renouvelable ; le charbon d’une ressource fossile.
\item Non : fabrication, usages et fin de vie peuvent avoir des impacts.
\end{questions}
\end{corrige}
\end{minipage}
\par\noindent\begin{minipage}{\linewidth}
\begin{corrige}{6}{Un bilan de lampe}
\begin{questions}
\item $200-60=140\unite J$.
\item $60+140=200\unite J$ : le bilan est conservatif.
\end{questions}
\par\smallskip\noindent\textbf{Pour comprendre.}\enspace Écrire le bilan avant de soustraire : $E_{\text{entrée}}=E_{\text{utile}}+E_{\text{autres}}$ sans stockage. Les autres transferts valent $200-60=140\unite J$. Contrôle : $60+140=200\unite J$.\par
\pcFigureEx{figures/cor-energie-lampe.pdf}{Sans stockage, la somme des transferts sortants égale le transfert entrant.}
\end{corrige}
\end{minipage}
\par\noindent\begin{minipage}{\linewidth}
\begin{corrige}{7}{Le moteur}
\begin{questions}
\item Oui : 350 + 150 = 500 J.
\item La mise en mouvement d’une charge, si tel est l’objectif du dispositif.
\end{questions}
\pcFigureEx{figures/cor-energie-moteur.pdf}{Le transfert mécanique peut être utile ; le transfert thermique ne disparaît pas du bilan.}
\end{corrige}
\end{minipage}
\par\noindent\begin{minipage}{\linewidth}
\begin{corrige}{8}{Un stockage augmente}
\begin{questions}
\item Le stockage augmente de $80-30=50\unite J$.
\item Non : cette égalité ne convient que lorsque le stockage ne varie pas.
\end{questions}
\par\smallskip\noindent\textbf{Pour comprendre.}\enspace La conservation s’écrit ici $\Delta E_{\text{stockée}}=E_{\text{entrée}}-E_{\text{sortie}}=80-30=50\unite J$. Ce n’est pas une disparition des 50 J : le système les stocke en plus.\par
\pcFigureEx{figures/cor-energie-stockage.pdf}{Entrée − sorties = variation de stockage : entrée et sortie ne sont pas égales si le stockage change.}
\end{corrige}
\end{minipage}
\par\noindent\begin{minipage}{\linewidth}
\begin{corrige}{9}{Une cellule photovoltaïque}
\begin{questions}
\item Elle reçoit de l’énergie par rayonnement et peut en transférer électriquement à un circuit.
\item Non : un rayonnement transporte aussi de l’énergie.
\end{questions}
\pcFigureEx{figures/cor-photovoltaique.pdf}{La lumière apporte de l’énergie par rayonnement ; la sortie électrique n’est pas la seule sortie réelle.}
\end{corrige}
\end{minipage}
\par\noindent\begin{minipage}{\linewidth}
\begin{corrige}{10}{Un bilan terrestre}
\begin{questions}
\item Son énergie stockée augmente.
\item Il faut connaître les valeurs et le périmètre des transferts ; des flèches qualitatives ne sont pas des mesures.
\end{questions}
\end{corrige}
\end{minipage}
\par\noindent\begin{minipage}{\linewidth}
\begin{corrige}{11}{Comparer deux lampes}
\begin{questions}
\item $100-40=60\unite J$, soit 60 \% de l’énergie de A.
\item L’éclairage utile, la durée et le périmètre de comparaison doivent être comparables.
\end{questions}
\end{corrige}
\end{minipage}
\par\noindent\begin{minipage}{\linewidth}
\begin{corrige}{12}{Une énergie cinétique fournie}
\begin{questions}
\item $E_c=\frac12\times2\times3^2=9\unite J$.
\item Non : avec $v=6\unite{m/s}$, on obtient $36\unite J$, soit quatre fois plus. La formule est fournie et constitue un approfondissement.
\end{questions}
\end{corrige}
\end{minipage}
\par\noindent\begin{minipage}{\linewidth}
\begin{corrige}{13}{Le mot « perte »}
\begin{questions}
\item Les transferts thermiques vers l’environnement conservent l’énergie, mais ne répondent pas au but de mise en mouvement.
\item Non : les formes et conditions de transfert ne sont pas également utilisables ; récupérer une énergie dispersée peut être difficile.
\end{questions}
\end{corrige}
\end{minipage}
\par\noindent\begin{minipage}{\linewidth}
\begin{corrige}{14}{Un raccourci sur le climat}
\begin{questions}
\item Ils absorbent et émettent une partie du rayonnement infrarouge ; ils ne bloquent pas toute sortie d’énergie.
\item Il faut comparer l’ensemble des transferts entrants et sortants pour le système Terre-atmosphère.
\end{questions}
\end{corrige}
\end{minipage}
\par\noindent\begin{minipage}{\linewidth}
\begin{corrige}{15}{Une comparaison incomplète}
\begin{questions}
\item Non : l’usage n’est qu’une partie du cycle de vie.
\item Extraction des matériaux, fabrication, transport, énergie fournie, durée de vie, réparabilité et traitement final.
\end{questions}
\end{corrige}
\end{minipage}
\par\noindent\begin{minipage}{\linewidth}
\phantomsection\label{cor-defi-12}\pcSolutionDefi{12}{Le convertisseur mystérieux}\begin{questions}
\item $800-500-220=80\unite J$ stockés en plus, si les transferts indiqués sont complets.
\item L’énergie n’a pas disparu ; le système a pu stocker 80 J ou un transfert a été oublié.
\item Examiner un échauffement interne, l’évolution du fonctionnement et les autres sorties ; le seul écart de bilan ne prouve pas un mécanisme.
\end{questions}
\pcFigureEx{figures/cor-convertisseur.pdf}{Le stockage de 80 J est déduit seulement si les transferts indiqués sont complets ; une sortie oubliée reste une autre hypothèse.}
\end{minipage}
\par\noindent\begin{minipage}{\linewidth}
\phantomsection\label{cor-activite-12}\pcSousTitre{Pistes pour l’investigation}
Pour chaque convertisseur, nommer le système, la source et les transferts entrant et sortant. Une lampe reçoit électriquement et cède notamment par rayonnement et thermiquement ; un moteur transfère aussi mécaniquement. Un échauffement indique une sortie ou un stockage à discuter. Sans valeurs complètes, présenter un bilan qualitatif plutôt qu’un partage chiffré inventé.\par\end{minipage}
